\scanpage{76}
En effet, \deleted{\unsure{quitte|Ceci}} on peut supposer \inserted{$S,T$ connexes réduits, donc $T$ un tore,} \deleted{déployés : faire une extension \inserted{finie} de $k$, \unsure{que|pour} $T$ et le tore \unsure{maximal} $S$ de $S$}
On peut supposer aussi $S$ extension d'une VA par un tore $S_0$. Quitte à faire une extension finie, on peut supposer
\[
S_0\simeq\mathbf G_m^r,\qquad T\simeq\mathbf G_m^s.
\]
Alors $T_\ell(S_0)\unsure{\times|<}T_\ell(T)\mathbin{-}T_\ell(\mathbf G_m)$ est un accouplement
\[
\mathbf Z_\ell^r\otimes\mathbf Z_\ell^s\mathbin{-}\mathbf Z_\ell^{\unsure{s|1}},
\]
où \deleted{$\unsure{\mathbf Z_\ell}$} l'élém. $f$ de Frob. opère sur \deleted{$\mathbf Z_\ell^{\unsure{s|r}}$} les trois facteurs par $q\,\mathrm{id},q\,\mathrm{id},q\,\mathrm{id}$, donc sur le premier \unsure{mbr|mbre|mb.} par $q^2\,\mathrm{id}$, sur le deuxième par $q\,\mathrm{id}$. Donc l'accouplement induit est nul. Cela nous ramène au cas où $S$ est une V.A. Mais par Weil, les val. propres de $f$ opérant sur $\deleted{\unsure{\mathbf Z_\ell|\mathbf Z|\mathbf Q}}T_\ell(S)$ sont de val. absolue $q^{1/2}$, donc les val. propres sur $T_\ell(S)\otimes T_\ell(\mathbf G_m)$ sont de val. absolue $q^{3/2}$, or sur $T_\ell(\mathbf G_m)$ on a $q$, on a encore \uuline{nullité}.

Nous traitons maintenant le cas général, \deleted{\unsure{Pour} \unsure{cela}}, qui, nous le savons, se ramène au cas $X$ un trait \inserted{$=\operatorname{Spec}V$, $k$ alg. clos} $Y$ son pt fermé. On écrit
\[
V=\underline{\lim} A_i,\qquad A_i\subset V,\qquad \unsure{A_i\ni\deleted{t\,\unsure{?}}|A_i\deleted{\text{étant}}}\text{ de t.f. sur }\mathbf Z.
\]
\deleted{\unsure{Pour} $i$ \unsure{grand},} Soit
\[
X_i=\unsure{\mathrm{Sp}|\operatorname{Spec}} A_i,\quad Y_i=V(\unsure{h_i|t_i|k_i|\mathfrak p_i}),\quad U_i=X_i-Y_i,
\]
donc
\[
X\unsure{=|\simeq}\varprojlim X_i,\quad Y=\varprojlim Y_i,\quad U=\varprojlim U_i.
\]
\margin{75}
